\subsection{商模的平坦性}

命题 4. 设 E 是右 A-模。下列三个性质等价：

\begin{enumerate}
\def\labelenumi{(\alph{enumi})}
\item
  E 是平坦的。
\item
  对形如
\end{enumerate}

\[
0 \longrightarrow \mathrm{G} \stackrel {v} {\longrightarrow} \mathrm{H} \stackrel {w} {\longrightarrow} \mathrm{E} \longrightarrow 0\tag{1}
\]

的右 A-模的每个正合列以及每个左 A-模 F，序列

\[
0 \longrightarrow \mathrm{G} \otimes_ {\mathrm{A}} \mathrm{F} \xrightarrow {v \otimes 1} \mathrm{H} \otimes_ {\mathrm{A}} \mathrm{F} \xrightarrow {w \otimes 1} \mathrm{E} \otimes_ {\mathrm{A}} \mathrm{F} \longrightarrow 0\tag{2}
\]

是正合的。

\begin{enumerate}
\def\labelenumi{(\alph{enumi})}
\setcounter{enumi}{2}
\tightlist
\item
  存在一个正合列 (1)，其中 H 是平坦的，使得对每个形如 \(A_{s}/a\) 的左 A-模 F，序列 (2) 都正合，这里 a 是 A 的一个有限生成左理想。
\end{enumerate}

我们先证明 (a) 蕴含 (b)。左 A-模 F 同构于一个自由模的商（《代数学》第二章 §1 第 11 目命题 20）；换句话说，我们有正合列

\[
0 \longrightarrow \mathrm{R} \stackrel {i} {\longrightarrow} \mathrm{L} \stackrel {p} {\longrightarrow} \mathrm{F} \longrightarrow 0
\]

其中 L 是自由的。考虑图

\[
\begin{array}{c} \mathrm{G} \otimes \mathrm{R} \xrightarrow {v \otimes 1 _ {\mathrm{R}}} \mathrm{H} \otimes \mathrm{R} \xrightarrow {w \otimes 1 _ {\mathrm{R}}} \mathrm{E} \otimes \mathrm{R} \\ \downarrow_ {1 _ {\mathrm{G}} \otimes i} \\ \mathrm{G} \otimes \mathrm{L} \xrightarrow {v \otimes 1 _ {\mathrm{L}}} \mathrm{H} \otimes \mathrm{L} \xrightarrow {w \otimes 1 _ {\mathrm{L}}} \mathrm{E} \otimes \mathrm{L} \\ \downarrow_ {1 _ {\mathrm{G}} \otimes p} \\ \mathrm{G} \otimes \mathrm{F} \xrightarrow {v \otimes 1 _ {\mathrm{F}}} \mathrm{H} \otimes \mathrm{F} \end{array}\tag{3}
\]

立即可得这个图是交换的，且由第 1 目的引理 1，它的行与列都正合；此外，由于 \(\mathbf{1}_{\mathbf{G}} \otimes p\) 和 \(\mathbf{1}_{\mathbf{H}} \otimes p\) 是满射的（第 1 目，引理 1），我们有 \(\mathbf{G} \otimes \mathbf{F} = \operatorname{Coker}(\mathbf{1}_{\mathbf{G}} \otimes i)\) ，

\[
\mathbf {H} \otimes \mathbf {F} = \operatorname{Coker} (\mathbf {1} _ {\mathbf {H}} \otimes i);
\]

\(w \otimes l_{R}\) 是满射的（第 1 目，引理 1）；最后，由于 L 自由因而是平坦的， \(v \otimes l_{L}\) 是单射的。于是可以应用蛇形图（§1 第 4 目，命题 2 (iii)）证明存在正合列

\[
\operatorname{Ker} \left(\mathrm{l} _ {\mathrm{H}} \otimes i\right) \longrightarrow \operatorname{Ker} \left(\mathrm{l} _ {\mathrm{E}} \otimes i\right) \xrightarrow {d} \mathrm{G} \otimes \mathrm{F} \xrightarrow {v \otimes \mathrm{l} _ {\mathrm{F}}} \mathrm{H} \otimes \mathrm{F}.\tag{4}
\]

在此条件下，若 E 是平坦的，则 \(1_{E} \otimes i\) 是单射的，换句话说 \(\operatorname{Ker}(1_{\mathbf{E}} \otimes i) = 0\) ，而正合列 (4) 表明 \(v \otimes 1_{F}\) 是单射的，因此序列 (2) 是正合的（考虑到第 1 目的引理 1）。

由于 (b) 显然蕴含 (c)，剩下要证明 (c) 蕴含 (a)。在 \(\mathbf{R} = \mathfrak{a}\) 、 \(\mathbf{L} = \mathbf{A}_s\) 、 \(\mathbf{F} = \mathbf{A}_s / \mathfrak{a}\) 的情形考虑图 (3)，并应用正合列 (4)。由假设 \(v \otimes 1_{\mathbb{F}}\) 是单射的，因此 \(\operatorname{Im}(d) = 0\) ；此外，由于 H 是平坦的，我们有 \(\operatorname{Ker}(1_{\mathbb{H}} \otimes i) = 0\) ；于是序列 (4) 的正合性蕴含 \(\operatorname{Ker}(1_{\mathbb{E}} \otimes i) = 0\) ，换句话说， \(1_{\mathbb{E}} \otimes i\) 是单射的，这就证明了 E 是平坦的（第 3 目，注记 1）。

命题 5. 设 \(0 \to E' \xrightarrow{v} E \xrightarrow{w} E'' \to 0\) 是右 A-模的一个正合列。设 \(E''\) 是平坦的。则为使 \(E\) 是平坦的，必要且充分的是 \(E'\) 是平坦的。

设 \(u: F' \to F\) 是左 A-模的一个单同态。考虑图

\[
\begin{array}{c} \mathrm{E} ^ {\prime} \otimes \mathrm{F} ^ {\prime} \xrightarrow {v \otimes 1 _ {\mathrm{F} ^ {\prime}}} \mathrm{E} \otimes \mathrm{F} ^ {\prime} \xrightarrow {w \otimes 1 _ {\mathrm{F} ^ {\prime}}} \mathrm{E} ^ {\prime \prime} \otimes \mathrm{F} ^ {\prime} \\ 1 _ {\mathrm{E} ^ {\prime}} \otimes u \Bigg \downarrow \qquad \qquad \qquad \qquad \qquad \qquad \qquad \qquad \qquad \qquad \qquad \qquad \qquad \qquad \qquad \qquad \qquad \qquad \qquad \qquad \qquad \qquad \qquad \qquad \qquad \qquad \qquad \qquad \qquad \qquad \qquad \qquad \qquad \qquad \\ \mathrm{E} ^ {\prime} \otimes \mathrm{F} \xrightarrow [ v \otimes 1 _ {\mathrm{F}} ]{} \mathrm{E} \otimes \mathrm{F} \xrightarrow [ w \otimes 1 _ {\mathrm{F}} ]{} \mathrm{E} ^ {\prime \prime} \otimes \mathrm{F} \end{array}
\]

它是交换的且它的行是正合的（第 1 目，引理 1）。由于 \(\mathbf{E}''\) 是平坦的，

\(1_{\mathbf{E}^{\prime \prime}} \otimes u\) 是单射的；此外，命题 4 证明 \(v \otimes 1_{\mathbf{F}^{\prime}}\) 和 \(v \otimes 1_{\mathbf{F}}\) 都是单射的。在此条件下，若 E 是平坦的，则 \(1_{\mathbf{E}} \otimes u\) 是单射的，因而

\[
\left(1 _ {\mathbf {E}} \otimes u\right) \circ (v \otimes 1 _ {\mathbf {F} ^ {\prime}}) = (v \otimes 1 _ {\mathbf {F}}) \circ \left(1 _ {\mathbf {E} ^ {\prime}} \otimes u\right);
\]

也是单射的；由此得 \(\mathbf{1}_{\mathbf{E}'} \otimes u\) 是单射的，从而 \(\mathbf{E}'\) 是平坦的。反之，若 \(\mathbf{E}'\) 是平坦的，则 \(\mathbf{1}_{\mathbf{E}'} \otimes u\) 是单射的；于是由 §1 第 4 目命题 2 的推论 1 得 \(\mathbf{1}_{\mathbf{E}} \otimes u\) 是单射的，因此 \(\mathbf{E}\) 是平坦的。

\textbf{注记}\par

\begin{enumerate}
\def\labelenumi{(\arabic{enumi})}
\item
  E 和 \(E'\) 都平坦而 \(E''\) 不平坦是可能的，Z-模的例子 \(E = Z\) 、 \(E' = nZ\) 、 \(E'' = Z/nZ\) （ \(n \geqslant 2\) ）就表明了这一点。
\item
  平坦模的子模不一定是平坦模（习题 3）。
\end{enumerate}

